\section{习题与解答}
\begin{exercise}\label{ex:m2horizon}
飞船从地球由静止出发，持续保持正向固有加速度$\alpha$。地球在$t_e>0$发出追赶飞船的光。求飞船接收该信号的固有时，并判断地球所有未来信号是否都能追上。
\Diagram{fig:m2horizonq}{恒固有加速世界线及追赶它的光线}{\draw[axis](-.2,0)--(4.5,0)node[right]{$\alpha x/c^2$};\draw[axis](0,0)--(0,4.6)node[above]{$\alpha t/c$};\draw[curve,domain=0:2,samples=90]plot({cosh(\x)-1},{sinh(\x)});\draw[->](0,.5)--(.25,.75);\draw[guide](0,1)--(3.4,4.4);\node[left]at(0,1){$1$};}
\end{exercise}
\begin{solution}
飞船世界线$ct=L_\alpha\sinh\chi$、$x=L_\alpha(\cosh\chi-1)$，$\chi=\alpha\tau/c$。图\ref{fig:m2horizonq}的信号满足$x=c(t-t_e)$，代入后
\[
 \frac{\alpha t_e}c=\sinh\chi-\cosh\chi+1=1-\ee^{-\chi}.
\]
所以仅当$0<t_e<c/\alpha$有有限接收固有时
\[
 \tau_R=-\frac c\alpha\ln(1-\alpha t_e/c).
\]
\Diagram{fig:m2horizona}{发射时刻逼近$c/\alpha$时，接收固有时发散}{\draw[axis](0,0)--(4.5,0)node[right]{$\alpha t_e/c$};\draw[axis](0,0)--(0,3.8)node[above]{$\alpha\tau_R/c$};\draw[curve,domain=0:.96,samples=90]plot({3.5*\x},{-ln(1-\x)});\draw[guide](3.5,0)--(3.5,3.5);\node[below]at(3.5,0){$1$};}
图\ref{fig:m2horizona}在边界发散。$t_e\geq c/\alpha$的信号永不追上持续加速的飞船，虽然飞船在每个有限时刻速度都小于$c$。这来自其未来世界线不断趋近一条光状渐近线。若飞船后来停机或转向，未来轨迹改变，不能继续使用永久加速的信号边界。
\end{solution}

\begin{exercise}\label{ex:m2labconstant}
光子火箭从静止出发，按初始实验室时间计算的后向光功率恒为$P$。求火箭剩余静质量和速度随实验室时间的变化，并与恒固有功率模型区分。初始质量$M_0$。
\Diagram{fig:m2labconstantq}{功率固定在同一实验室中定义}{\draw[curve](2,0)rectangle(4,1);\draw[->](1.8,.5)--(.1,.5)node[above]{$P$实验室};\draw[->](4.2,.5)--(5.7,.5)node[above]{$v(t)$};}
\end{exercise}
\begin{solution}
图\ref{fig:m2labconstantq}截至时间$t$发出的总光能$Pt$、总后向动量$-Pt/c$。全系统守恒给飞船总能量$E=M_0c^2-Pt$、动量$p=Pt/c$。令$s=Pt/(M_0c^2)$，质壳关系得到
\[
 m=M_0\sqrt{1-2s},\qquad \beta=\frac{pc}E=\frac s{1-s},\quad 0\leq s<1/2.
\]
实际还须在$m$达到干质量时停机，不能把$m=0$极限当有质量飞船状态。
\Diagram{fig:m2labconstanta}{实验室定功率与固有定功率对应不同质量时间关系}{\draw[axis](0,0)--(4.4,0)node[right]{$s$};\draw[axis](0,0)--(0,2.8)node[above]{$m/M_0$};\draw[curve,domain=0:.499,samples=90]plot({7*\x},{2.3*sqrt(1-2*\x)});\node[below]at(3.5,0){$1/2$};\node[left]at(0,2.3){$1$};}
图\ref{fig:m2labconstanta}的质量不是实验室时间的线性函数。由$P=P_*(1-\beta)$，所需固有功率$P_*=P(1-s)/(1-2s)$会增大。固有时间满足$\dd\tau=\sqrt{1-2s}\dd t/(1-s)$。令$K=M_0c^2/P$、$w=\sqrt{1-2s}$，则$\dd t=K\dd s=-Kw\dd w$、$1-s=(1+w^2)/2$，所以
\[
 \tau=-2K\int_1^w\frac{q^2}{1+q^2}\dd q
 =-2K[q-\arctan q]_1^w.
\]
代入上下限，得到
\[
 \tau=\frac{2M_0c^2}P\left[1-w+\arctan w-\frac\pi4\right].
\]
在$t=0$有$w=1,\tau=0$。可见恒功率条件随参考系而变。固有功率恒定时采用$m=M_0-P_*\tau/c^2$，对应另一种推进程序。
\end{solution}

\begin{exercise}\label{ex:m2midburn}
光子火箭需从静止到达距离$L$处并停止，前半程保持固有加速度$\alpha$，后半程保持同样大小的固有减速度。求地球时间、船上时间和初末质量比。
\Diagram{fig:m2midq}{中点反转推力，速度连续}{\draw[axis](-.3,0)--(5.8,0)node[right]{$x$};\fill(0,0)circle(1.5pt);\fill(5,0)circle(1.5pt);\draw[guide](2.5,-.7)--(2.5,1.3);\draw[->](.5,.6)--(1.8,.6)node[above]{$\alpha$};\draw[->](4.5,.6)--(3.2,.6)node[above]{$-\alpha$};}
\end{exercise}
\begin{solution}
对图\ref{fig:m2midq}的前半程，用恒固有加速度世界线，令$q=\alpha L/(2c^2)$，得到中点快度$\chi_m=\operatorname{arcosh}(1+q)$。该段地球时间$c\sinh\chi_m/\alpha$、船上时间$c\chi_m/\alpha$。
\Diagram{fig:m2mida}{固有时中的快度先线性增加再线性减少}{\draw[axis](0,0)--(5.2,0)node[right]{$\tau$};\draw[axis](0,0)--(0,3)node[above]{$\chi$};\draw[curve](0,0)--(2.3,2.5)--(4.6,0);\node[above]at(2.3,2.5){$\chi_m$};}
后半程是时间反向的减速过程，图\ref{fig:m2mida}两段持续时间相同。因此$T=2c\sqrt{q(2+q)}/\alpha$、$\tau=2c\operatorname{arcosh}(1+q)/\alpha$。光子火箭消耗与总绝对快度$2\chi_m$对应，$M_0/M_f=\ee^{2\chi_m}=[1+q+\sqrt{q(2+q)}]^2$。减速阶段仍耗质量，不能让快度符号变化抵消燃料消耗。
\end{solution}

\begin{exercise}\label{ex:m2staging}
固定排气速度$u$的火箭初质量100，燃烧至60时丢弃质量20的空壳，再燃烧至最终质量20。求总快度，并与同样初质量、最终携带40质量但不丢壳的单级模型比较。所有质量使用同一单位。
\Diagram{fig:m2stageq}{无相对弹射的丢壳瞬间，飞船速度不变}{\node at(0,0){$100$};\draw[->](.6,0)--(1.5,0);\node at(2,0){$60$};\draw[->](2.6,0)--(3.5,0);\node at(4,0){$40$};\draw[->](4.6,0)--(5.5,0);\node at(6,0){$20$};\node[below]at(3,-.6){丢壳20};}
\end{exercise}
\begin{solution}
图\ref{fig:m2stageq}第一段快度$(u/c)\ln(100/60)$。丢壳时空壳与飞船相对静止，四动量按共同四速度分开，没有快度突变；第二段快度$(u/c)\ln(40/20)$。总快度$\chi=(u/c)\ln(10/3)$。
\Diagram{fig:m2stagea}{对数质量图中竖直丢壳不贡献推进快度}{\draw[axis](0,0)--(5.2,0)node[right]{阶段};\draw[axis](0,0)--(0,3.5)node[above]{$\ln m$};\draw[curve](0,3)--(1.8,2.3)--(1.8,1.7)--(4.2,.7);\node[right]at(1.8,2){丢壳};}
不丢壳的对照模型结束时含壳和有效载荷总质量40。其快度为
\[
 \chi_1=\frac uc\ln\frac{100}{40},\qquad
 \chi-\chi_1=\frac uc\ln\frac43>0.
\]
图\ref{fig:m2stagea}对应的分级增益为后一式。若把无推进丢掉的20也按火箭式积分，会错误地多算一段快度；比较模型的有效载荷、燃料和空壳质量必须逐项一致。
\end{solution}

\begin{exercise}\label{ex:m2labpower}
光子火箭沿$+x$运动，瞬时静止系向后辐射功率为$P_*$。求按实验室发射时间计算的光功率，并由它和飞船动能变化核对质量能守恒。
\Diagram{fig:m2labpowerq}{后向光子在实验室红移}{\draw[curve](2,0)rectangle(4,1);\draw[axis](4.2,.5)--(5.6,.5)node[above]{$V$};\draw[axis](1.8,.5)--(.2,.5)node[above]{$P_*$};}
\end{exercise}
\begin{solution}
图\ref{fig:m2labpowerq}后向光子能量因子为$\gamma(1-\beta)$，发射时间因子为$\gamma$，所以$P_\gamma=P_*(1-\beta)$。同时$\dd m/\dd t=-P_*/(\gamma c^2)$，飞船快度率$\dd\chi/\dd t=P_*/(\gamma mc^2)$。
\Diagram{fig:m2labpowera}{飞船总能量减小率等于后向辐射实验室功率}{\draw[curve](0,0)rectangle(2.2,1);\node at(1.1,.5){$\gamma mc^2$};\draw[->](2.4,.5)--(3.5,.5);\node[right]at(3.7,.5){$P_*(1-\beta)$};}
对图\ref{fig:m2labpowera}的飞船总能量求导，利用$\dd\gamma/\dd\chi=\gamma\beta$，有
\[
 \frac{\dd(\gamma mc^2)}{\dd t}
 =c^2\left(\gamma\frac{\dd m}{\dd t}
 +m\frac{\dd\gamma}{\dd\chi}\frac{\dd\chi}{\dd t}\right)
 =-P_*+mc^2\gamma\beta\frac{P_*}{\gamma mc^2}
 =-P_*(1-\beta).
\]
飞船虽加速，总能量仍因辐射流失而下降；静质量损失的一部分转为更大的运动因子。不能把“速度增加”直接等同于飞船总能量必增。
\end{solution}

\begin{exercise}\label{ex:m2relay}
沿用地球与恒星相距10光年、速度$0.8c$的过程，但旅行者不转向。另一只返航钟在恒星处与其相遇，把读数设置为7.5年后飞回地球。比较重逢读数与原双生子结果，解释是否必须由同一只钟经历加速才能得到差值。
\Diagram{fig:m2relayq}{两只匀速钟在恒星处交接读数}{\draw[axis](-.2,0)--(4.2,0);\draw[axis](0,0)--(0,5.4)node[above]{$ct$};\draw[curve](0,0)--(2,2.5)--(3.7,4.625);\draw[curve](3.6,.5)--(0,5);\fill(2,2.5)circle(1.4pt)node[right]{$B$交接};}
\end{exercise}
\begin{solution}
图\ref{fig:m2relayq}去程钟从出发到$B$经历7.5年，返航钟从$B$到地球也经历7.5年。它在$B$被设置为7.5年，所以到地球读15年，地球仍为25年。
\Diagram{fig:m2relaya}{交接后的读数拼接两段不同世界线的固有时间}{\draw[axis](0,0)--(5.1,0)node[right]{地球时间};\draw[axis](0,0)--(0,3)node[above]{拼接读数};\draw[curve](0,0)--(4.6,2.4);\draw[guide](2.3,0)--(2.3,1.2);\node[below]at(2.3,0){12.5年};\node[right]at(4.6,2.4){15年};}
图\ref{fig:m2relaya}的拼接读数沿两段组成的折线累积，几何上与瞬时转向世界线相同。没有单只返航钟需要在$B$改变速度，但必须实际校准读数；两只钟若早在各自惯性系随意同步，初始偏移未必满足本题。差值由选择的时空路径决定，加速是连接路径的一种实现方式。
\end{solution}

\begin{exercise}\label{ex:m2retardedzero}
发光点在实验室沿$(vt,b)$运动，观察者在原点。观察者在$t_o=0$拍照，而该坐标时刻物体恰在最近点$x=0$。求照片中光的发射时刻和发射位置。
\Diagram{fig:m2zeroq}{同时最近点与照片发光点不是同一点}{\draw[axis](-3,0)--(2.5,0)node[right]{$x$};\draw[axis](0,0)--(0,2.8)node[above]{$y$};\draw[curve](-2.6,1.6)--(2.2,1.6);\fill(0,1.6)circle(1.5pt)node[above right]{同时位置};}
\end{exercise}
\begin{solution}
图\ref{fig:m2zeroq}的延迟条件为$0=t_e+\sqrt{v^2t_e^2+b^2}/c$，要求$t_e<0$。移项平方后$c^2t_e^2=v^2t_e^2+b^2$，所以$t_e=-\gamma b/c$、$x_e=-\beta\gamma b$。
\Diagram{fig:m2zeroa}{入眼光线来自最近点之前，路程为$\gamma b$}{\draw[axis](-3.2,0)--(2.3,0)node[right]{$x$};\draw[axis](0,0)--(0,2.7)node[above]{$y$};\fill(-2,1.6)circle(1.5pt)node[above]{$x_e=-\beta\gamma b$};\draw[->](-2,1.6)--(0,0);\draw[guide](0,0)--(0,1.6);\node[right]at(0,.8){$b$};}
代回原式，光程$R_e=\sqrt{\beta^2\gamma^2b^2+b^2}=\gamma b$，恰在$\gamma b/c$后到达，验证选取了正确负根。图\ref{fig:m2zeroa}中的源虽在坐标时间0已到最近点，照片仍显示此前位置；照片记录的是经过传播后到达的旧信号。
\end{solution}

\begin{exercise}\label{ex:m2fuelbudget}
理想光子火箭总初末质量比为$R$，需把给定质量的有效载荷送到远方并相对地球停下。不规定距离，忽略巡航损失。求最大巡航洛伦兹因子以及质量比达到100时的巡航速度。
\Diagram{fig:m2budgetq}{燃料预算分成加速和减速两个相等快度段}{\draw[axis](0,0)--(5.2,0);\draw[curve](0,.3)--(2.3,1.8)--(4.6,.3);\node[above]at(2.3,1.8){巡航};\node[below]at(0,0){静止};\node[below]at(4.6,0){静止};}
\end{exercise}
\begin{solution}
图\ref{fig:m2budgetq}总绝对快度为$\ln R$，要最终停止，两段均为$\chi=\tfrac12\ln R$。所以$\gamma_{\max}=\cosh(\tfrac12\ln R)=(\sqrt R+1/\sqrt R)/2$，$\beta=(R-1)/(R+1)$。
\Diagram{fig:m2budgeta}{保留制动燃料会显著降低同一初末质量比的巡航速度}{\draw[axis](0,0)--(4.8,0)node[right]{$R$};\draw[axis](0,0)--(0,2.8)node[above]{$\beta$};\draw[curve,domain=1:100,samples=80]plot({.04*\x},{2.3*(\x-1)/(\x+1)});\node[below]at(4,0){100};}
取$R=100$，$\gamma=5.05$、$\beta=99/101\simeq.980198$，如图\ref{fig:m2budgeta}。若把全部燃料只用来加速，速度可达$(R^2-1)/(R^2+1)$，但剩余质量不足以再完成制动。任务若要求到达后停下，质量比还要包含制动段的消耗。
\end{solution}

\begin{exercise}\label{ex:m2signals}
例\ref{ex:m2twins}中地球每年发出一束时间标记光信号。求旅行者去程和返程接收到的频率，并据此还原地球25年的连续钟信号。
\Diagram{fig:m2signalsq}{地球连续发光，旅行者先远离后靠近}{\draw[axis](-.2,0)--(3.5,0)node[right]{$x$};\draw[axis](0,-.2)--(0,5.5)node[above]{$ct$};\draw[curve](0,0)--(2,2.5)--(0,5);\foreach\y in{.1,.3,.5}{\draw[guide](0,\y)--({4*\y},{5*\y});}\foreach\y in{1.3,2.1,2.9,3.7,4.5}{\draw[guide](0,\y)--({(4-.8*\y)/1.8},{(4+\y)/1.8});}}
\end{exercise}
\begin{solution}
去程接收者远离，频率比为$\sqrt{(1-\beta)/(1+\beta)}=1/3$；返程靠近，频率比为3。旅行者两段固有时间各7.5年，所以去程接收相当于地球2.5年内发出的信号，返程接收其余22.5年。
\Diagram{fig:m2signalsa}{转向瞬间接收到的光来自地球2.5年时刻}{\draw[axis](-.2,0)--(3.5,0)node[right]{$x$};\draw[axis](0,-.2)--(0,5.5)node[above]{$ct$};\draw[curve](0,0)--(2,2.5)--(0,5);\draw[very thick,->](0,.5)--(2,2.5);\node[left]at(0,.5){2.5年};\node[right]at(2,2.5){$B$};}
图\ref{fig:m2signalsa}中由光速连线也可算发射时刻$t_e=t_B-l/c=12.5-10=2.5$年。与旅行者去程、返程赋予地球的同时刻4.5年、20.5年都不同。接收信号连续，只有接收频率在理想瞬时换向模型中突变。
\end{solution}

\begin{exercise}\label{ex:m2fuelreturn}
光子火箭需由地球静止起飞，加速至$\beta c$，到远处减速至相对地球静止，随后原路返回并最终停下。忽略巡航耗能，求包括全部返航燃料在内的初末质量比。
\Diagram{fig:m2returnq}{四段速度改变与两段巡航}{\draw[axis](-.3,0)--(6,0)node[right]{$x$};\fill(0,0)circle(1.5pt)node[below]{地球};\fill(5,0)circle(1.5pt)node[below]{远端};\draw[->](.3,.8)--(4.7,.8)node[midway,above]{$+\beta c$};\draw[->](4.7,-.8)--(.3,-.8)node[midway,below]{$-\beta c$};}
\end{exercise}
\begin{solution}
令$\chi=\operatorname{artanh}\beta>0$。图\ref{fig:m2returnq}中每次由静止到巡航或由巡航到静止，都需消耗绝对快度变化$\chi$。喷流方向在减速时反转，火箭质量仍减少，故应对快度变化取绝对值。
\Diagram{fig:m2returna}{快度变化总长度为$4\chi$}{\draw[axis](-.3,0)--(6,0)node[right]{过程};\draw[axis](0,-1.5)--(0,1.8)node[above]{$\chi_{\rm ship}$};\draw[very thick](0,0)--(1,1)--(2,1)--(3,0)--(4,-1)--(5,-1)--(6,0);\node[left]at(0,1){$\chi$};\node[left]at(0,-1){$-\chi$};}
图\ref{fig:m2returna}中四次点火的质量比逐次相乘，故$M_0/M_f=\ee^{4\chi}=[(1+\beta)/(1-\beta)]^2$。若$\beta=.8$，质量比为81。不能先算单程消耗，再以最终干质量为基准简单加倍，因为起飞时必须推动全部后续燃料。
\end{solution}

\begin{exercise}\label{ex:m2constalpha}
飞船从静止出发，沿$x$保持固有加速度$\alpha$，直到到达地球距离$L$处。求地球时间、船上时间和末速度。
\Diagram{fig:m2alphaq}{固定终点距离的持续加速}{\draw[axis](-.3,0)--(5,0)node[right]{$x$};\fill(0,0)circle(1.5pt);\draw[curve](4,-.5)--(4,1.2);\draw[<->](0,-.5)--(4,-.5)node[midway,below]{$L$};\draw[axis](.5,.5)--(2,.5)node[above]{$\alpha$};}
\end{exercise}
\begin{solution}
图\ref{fig:m2alphaq}取飞船由原点静止出发。固有加速度定义给$\dd\chi/\dd\tau=\alpha/c$，所以$\chi=\alpha\tau/c$。再由四速度的时间和空间分量，
\[
 t=\int_0^\tau\cosh(\alpha s/c)\dd s
 =\left[\frac c\alpha\sinh(\alpha s/c)\right]_0^\tau
 =\frac c\alpha\sinh\chi,
\]
\[
 x=\int_0^\tau c\sinh(\alpha s/c)\dd s
 =\left[\frac{c^2}\alpha\cosh(\alpha s/c)\right]_0^\tau
 =\frac{c^2}\alpha(\cosh\chi-1).
\]
两个积分常数已由$t=x=0$确定。
\Diagram{fig:m2alphaa}{终点落在恒固有加速度双曲线上}{\draw[axis](-.2,0)--(4.6,0)node[right]{$\alpha x/c^2$};\draw[axis](0,-.2)--(0,4.6)node[above]{$\alpha t/c$};\draw[curve,domain=0:2,samples=70]plot({cosh(\x)-1},{sinh(\x)});\fill(2,2.828)circle(1.5pt)node[right]{$x=L$};\draw[guide](2,0)--(2,2.828);}
在图\ref{fig:m2alphaa}中令$q=\alpha L/c^2$。到达时$\cosh\chi_f=1+q$，正向运动选择$\sinh\chi_f=+\sqrt{(1+q)^2-1}=\sqrt{q(2+q)}$，因此
\[
 t_f=\frac c\alpha\sqrt{q(2+q)},\quad
 \tau_f=\frac c\alpha\operatorname{arcosh}(1+q),\quad
 \beta_f=\frac{\sqrt{q(2+q)}}{1+q}.
\]
平方后$1-\beta_f^2=(1+q)^{-2}>0$，也给$\gamma_f=1+q$。小$q$时$\operatorname{arcosh}(1+q)\simeq\sqrt{2q}$，两种时间同趋$\sqrt{2L/\alpha}$、末速趋$\sqrt{2\alpha L}$。大$q$时$t_f\simeq L/c+c/\alpha$，而$\tau_f\simeq(c/\alpha)\ln(2q)$，远距离的船上时间仅呈对数增长；这依赖全程持续提供恒固有加速度。
\end{solution}

\begin{exercise}\label{ex:m2powerengine}
理想光子火箭在自身瞬时静止系的辐射功率恒为$P_*$。将它与恒固有加速度光子火箭比较：分别求质量随固有时间的变化及固有功率随时间的条件。
\Diagram{fig:m2engineq}{定向光束的固有功率与推力}{\draw[curve](2,0)rectangle(4.3,1);\node at(3.15,.5){火箭$m$};\draw[axis](1.8,.5)--(.2,.5)node[above]{$P_*/c$};\draw[axis](4.5,.5)--(5.7,.5)node[above]{$\alpha$};}
\end{exercise}
\begin{solution}
图\ref{fig:m2engineq}中，瞬时船系向后发光的能量率为$P_*$，动量率大小为$P_*/c$。一阶能量与动量守恒分别给
\[
 \frac{\dd m}{\dd\tau}=-\frac{P_*}{c^2},\quad
 mc\frac{\dd\chi}{\dd\tau}=\frac{P_*}{c},\quad
 \alpha=c\frac{\dd\chi}{\dd\tau}=\frac{P_*}{mc}.
\]
若$P_*$恒定，以$m(0)=M_0$积分：$\int_{M_0}^{m}\dd\mu=-(P_*/c^2)\int_0^\tau\dd s$，得到$m=M_0-P_*\tau/c^2$。再相除两个微分方程，有$\dd\chi=-\dd m/m$，从初始静止积分得$\chi=-[\ln\mu]_{M_0}^{m}=\ln(M_0/m)$。推力恒定而惯性质量减小，故固有加速度不断增大。
\Diagram{fig:m2enginea}{恒加速度需要功率与剩余质量同步下降}{\draw[axis](0,0)--(4.3,0)node[right]{$\alpha\tau/c$};\draw[axis](0,0)--(0,2.8)node[above]{$m/M_0,\ P_*/P_{*0}$};\draw[curve,domain=0:3.5,samples=70]plot(\x,{2.3*exp(-\x)});\node[left]at(0,2.3){$1$};}
若要求图\ref{fig:m2enginea}中的$\alpha$恒定，必须令功率随质量改变：$P_*=mc\alpha$。质量方程于是变为
\[
 \frac{\dd m}{m}=-\frac\alpha c\dd\tau,
 \quad \int_{M_0}^{m}\frac{\dd\mu}{\mu}
 =-\frac\alpha c\int_0^\tau\dd s,
 \quad \ln\frac m{M_0}=-\frac{\alpha\tau}c.
\]
取指数并代回功率条件，
\[
 m=M_0\ee^{-\alpha\tau/c},\qquad
 P_*=M_0c\alpha\ee^{-\alpha\tau/c}.
\]
此时$\chi=\alpha\tau/c$，与恒固有加速度世界线一致。若干质量为$M_d$，两种程序都在$m=M_d$时停机：恒功率的可燃固有时间为$(M_0-M_d)c^2/P_*$，恒加速度则为$(c/\alpha)\ln(M_0/M_d)$。在相同初末质量比下，两种程序给出相同末快度；燃烧历时和沿途位置由功率程序决定。
\end{solution}

\begin{exercise}\label{ex:m2oppositejets}
一对同时由同一点射出的对称光斑在源系以同速$\beta c$沿相反方向运动。视线与靠近光斑速度夹角为$\theta$。测得两侧视横向速度分别为$u_+,u_-$。求$\beta\cos\theta$及$\beta\sin\theta$。
\Diagram{fig:m2jetsq}{靠近和远离光斑的投影}{\draw[axis](-.3,0)--(4.5,0)node[right]{视线};\draw[very thick,->](0,0)--(2.7,1.7)node[right]{$+$};\draw[very thick,->](0,0)--(-1.9,-1.2)node[left]{$-$};\node[below right]at(0,0){源};}
\end{exercise}
\begin{solution}
图\ref{fig:m2jetsq}的两侧横向位移大小相同，但接收时间分别缩短和拉长，故
\[
 \frac{u_+}c=\frac{\beta\sin\theta}{1-\beta\cos\theta},\quad
 \frac{u_-}c=\frac{\beta\sin\theta}{1+\beta\cos\theta}.
\]
当$u_++u_->0$时，两式相除后解得$A=\beta\cos\theta=(u_+-u_-)/(u_++u_-)$。再乘回任一分母，得$B=\beta\sin\theta=2u_+u_-/[c(u_++u_-)]$。若两侧视横向速度都为零，可能是$\beta=0$，也可能是运动轴沿视线；此时横向数据不足以唯一反演，不能使用上述零分母公式。
\Diagram{fig:m2jetsa}{两个可测投影组合给出速度和视角}{\draw[axis](0,0)--(3.8,0)node[right]{$A$};\draw[axis](0,0)--(0,2.8)node[above]{$B$};\draw[very thick,->](0,0)--(2.7,1.8)node[right]{$\beta$};\draw[guide](2.7,0)--(2.7,1.8)--(0,1.8);\node at(.9,.3){$\theta$};}
图\ref{fig:m2jetsa}给$\beta=\sqrt{A^2+B^2}$、$\theta=\operatorname{atan2}(B,A)$。若由观测反演得到$\beta\geq1$，则对称、同速、远场或距离标定中至少一项不满足模型假设，不能把不相容的数据解释成真实超光速。
\end{solution}

\begin{exercise}\label{ex:m2imagebar}
长度$L_0$的细杆沿$x$运动，杆与速度平行。观察者在远处，视线单位矢量$\vct n$与$x$夹角$\theta$。求同一照片中杆端横向投影的长度，并与等时长度比较。
\Diagram{fig:m2barq}{照片由两端不同的发光事件组成}{\draw[curve](0,0)--(3.5,0);\fill(0,0)circle(1.5pt)node[below]{$A$};\fill(3.5,0)circle(1.5pt)node[below]{$B$};\draw[axis](1.7,.5)--(3,.5)node[above]{$V$};\draw[guide](0,0)--(1.5,2.2);\draw[guide](3.5,0)--(5,2.2);\node[right]at(5,2.2){$\vct n$};}
\end{exercise}
\begin{solution}
图\ref{fig:m2barq}中，实验室同时杆长为$L=L_0/\gamma$。取尾端、头端世界线$x_A=Vt_A$、$x_B=Vt_B+L$。远处观察者的方向$\vct n$由物体指向相机，源点到相机的距离差近似为$-\vct n\cdot\Delta\vct r_e=-\cos\theta\Delta x_e$。同一照片要求
\[
 0=\Delta t_o=t_B-t_A-\frac{\cos\theta}{c}\Delta x_e.
\]
再由两端世界线相减，$\Delta x_e=V(t_B-t_A)+L$，消去发射时间差得
\[
 \Delta x_e=\frac{L}{1-\beta\cos\theta}.
\]
\Diagram{fig:m2bara}{发光事件空间差向成像面的投影}{\draw[curve](0,0)--(4,0);\draw[guide](0,0)--(2,2.5);\draw[guide](4,0)--(5.56,1.95);\draw[very thick](2,2.5)--(4.44,.55);\draw[guide](4,0)--(4.44,.55);\node[above right]at(3.4,1.6){$L_{\rm photo}$};\node[below]at(2,0){$\Delta x_e$};}
在图\ref{fig:m2bara}的像平面内，端点分离矢量须扣去沿视线分量，故
\begin{align*}
 L_{\rm photo}^2
 &=|\Delta\vct r_e|^2-(\vct n\cdot\Delta\vct r_e)^2\\
 &=(\Delta x_e)^2\sin^2\theta,
\end{align*}
因此$L_{\rm photo}=L_0\sin\theta/[\gamma(1-\beta\cos\theta)]$。在$\theta=\pi/2$时恢复$L_0/\gamma$。其角度函数与例\ref{ex:m2super}相同，沿用该例已经展开的商式求导，最大值在$\cos\theta=\beta$取得；代入$\sin\theta=1/\gamma$，得$L_{\rm photo,max}=L_0$。式中的$L$、$\Delta x_e$、$L_{\rm photo}$分别表示同时长度、延迟发射位置差和像面投影长，三者须按各自定义使用。
\end{solution}
